\documentclass[11pt,a4paper]{article}
\usepackage[T1]{fontenc}
\usepackage[utf8]{inputenc}
\usepackage{lmodern}
\usepackage[margin=27mm]{geometry}
\usepackage{amsmath,amssymb,amsthm}
\usepackage{microtype}
\usepackage[numbers,sort&compress]{natbib}
\usepackage[colorlinks=true,linkcolor=blue,citecolor=blue,urlcolor=blue,filecolor=blue,anchorcolor=blue,menucolor=blue,breaklinks=true,pdftitle={The log-concavity to log-convexity transition in fixed columns of Stirling numbers},pdfauthor={Xiaojian Zeng}]{hyperref}
\newtheorem{theorem}{Theorem}
\newtheorem{lemma}[theorem]{Lemma}
\newtheorem{proposition}[theorem]{Proposition}
\theoremstyle{remark}
\newtheorem{remark}[theorem]{Remark}
\newcommand{\D}{D_k}
\title{The log-concavity to log-convexity transition\protect\\
in fixed columns of Stirling numbers}
\author{Xiaojian Zeng\thanks{PBC School of Finance, Tsinghua University,
100190, China. Email: \texttt{xiaojianzeng@foxmail.com}.}}
\date{}

\begin{document}
\maketitle
\begin{abstract}
Let $c(n,k)$ be the number of permutations of $n$ elements with $k$ cycles.
We prove that each fixed column changes from log-concavity to log-convexity,
as conjectured by Liang and Sagan. More precisely, for every $k\geq1$ there
is a finite $N_k\geq k+1$ such that
$c(n,k)^2-c(n-1,k)c(n+1,k)$ is nonnegative for $1\leq n<N_k$ and strictly
negative for $n\geq N_k$. The proof establishes discrete convexity in $n$
of $c(n,k-1)/c(n,k)$, strictly so for $k\geq2$. A preservation identity
for elementary-symmetric coefficient ratios and a positive decomposition
of the resulting numerator give this convexity. Monotonicity and telescoping
then produce an exact criterion for a finite transition index.
\end{abstract}

\section{Statement and context}

We use the signless Stirling numbers of the first kind, with
$c(0,0)=1$, $c(n,0)=0$ for $n>0$, and $c(n,k)=0$ when $k<0$ or $k>n$.
For $k\geq1$ and $n\geq1$, put
\[
  \D(n)=c(n,k)^2-c(n-1,k)c(n+1,k).
\]
Thus the sign of $\D(n)$ describes log-concavity or log-convexity at the
center $n$. Liang and Sagan conjectured that, for every fixed $k$, the
nonnegative signs precede the nonpositive signs
\cite{LiangSagan2026} (Conjecture~8.8, p.~31). Their published statement
reports verification for $1\leq k\leq n\leq100$; the earlier arXiv version
numbers the same assertion Conjecture~8.9. We prove the assertion on its
whole domain, with a strictly negative tail.

\begin{theorem}\label{thm:main}
For every integer $k\geq1$, define, for integers $n\geq k$,
\[
  q_k(n)=\frac{c(n,k-1)}{c(n,k)}.
\]
Then
\begin{equation}\label{eq:qconvex}
 q_k(n)-2q_k(n+1)+q_k(n+2)\geq0 \qquad(n\geq k),
\end{equation}
with strict inequality if $k\geq2$ and equality if $k=1$.
The integer
\begin{equation}\label{eq:threshold}
 N_k=1+\min\{j\geq k:q_k(j)-q_k(j+1)<1\}
\end{equation}
is well defined and satisfies $N_k\geq k+1$. Moreover,
\[
 \D(n)\geq0 \quad(1\leq n<N_k),\qquad
 \D(n)<0 \quad(n\geq N_k).
\]
In particular, $N_1=2$.
\end{theorem}

The distinction between rows and columns is essential. Salas proves
coefficientwise strong log-concavity of fixed rows of the
Graham--Knuth--Patashnik array and strong log-convexity of its sequence of
row-generating polynomials \cite{Salas2026} (Theorems~1.4 and~1.5).
For row polynomials $P_n(z)=\sum_j c(n,j)z^j$, the coefficient of $z^{2k}$
in $P_{n-1}(z)P_{n+1}(z)-P_n(z)^2$ is a sum over pairs of indices
$i,2k-i$. It does not isolate $-\D(n)$. The argument below supplies the
passage from inequalities within a row to convexity of $q_k(n)$ as the row
index varies.

Fan and Wang's abstract reports convexity of elementary-symmetric ratios
obtained through operators preserving real-rootedness \cite{FanWang2026}.
We give the precise ratio lemma needed here with its full proof and claim
no priority for that general property. Their complete theorem and possible
applications to Stirling numbers have not been compared, because the full
text was unavailable. The contribution presented here is the explicit
positive decomposition leading to \eqref{eq:qconvex} and the finite-threshold
argument for the stated Liang--Sagan conjecture; no worldwide priority
claim is made.

\section{Convexity of elementary-symmetric coefficient ratios}

For positive real numbers $x_1,\ldots,x_m$, let $e_j=e_j(x_1,\ldots,x_m)$
be defined by
\[
 \prod_{i=1}^m(1+x_i z)=\sum_{j=0}^m e_jz^j.
\]
We set $e_0=1$ and $e_j=0$ for $j<0$. A sequence $(R_j)_{j=0}^m$ is
discretely convex if $R_{j+1}-2R_j+R_{j-1}\geq0$ for $1\leq j<m$.

\begin{lemma}\label{lem:ratios}
For every $m\geq1$ and every positive $x_1,\ldots,x_m$, the sequence
\[
 R_0=0,\qquad R_j=\frac{e_{j-1}}{e_j}\quad(1\leq j\leq m)
\]
is strictly increasing and discretely convex.
\end{lemma}

\begin{proof}
For $t>0$, set
\[
 F_t(u,v)=\frac{u(1+tv)}{1+tu}.
\]
Suppose $0\leq u\leq v$, $h\geq2u-v$, and $w\geq2v-u$.
The function $F_t(w,v)$ increases with $w\geq0$, since its derivative
is $(1+tv)/(1+tw)^2>0$. Also $F_t(u,h)$ is nondecreasing in $h$ on
the entire real line, since its derivative is $tu/(1+tu)\geq0$.
Consequently,
\begin{align}
 &F_t(w,v)-2F_t(v,u)+F_t(u,h)\notag\\
 &\quad\geq F_t(2v-u,v)-2F_t(v,u)+F_t(u,2u-v)\notag\\
 &\quad=\frac{2t^2(v-u)^3}
 {(1+tu)(1+tv)(1+t(2v-u))}\geq0.\label{eq:preserve}
\end{align}
The equality follows by clearing the three positive denominators.
The comparison value $2u-v$ may be negative, but it occurs only as
the second argument of $F_t$, hence only in a numerator. The inequality
also holds for $w=+\infty$ if
\begin{equation}\label{eq:endpoint}
 F_t(+\infty,v)=v+1/t,
\end{equation}
by taking the limit in the first term.

We show that multiplication by $1+tz$ preserves the asserted properties.
Let $a_0,\ldots,a_m>0$, $m\geq1$, and assume that
$R_0=0$, $R_j=a_{j-1}/a_j$ form an increasing, convex sequence.
Put $a_{-1}=a_{m+1}=0$ and
$b_j=a_j+ta_{j-1}$ for $0\leq j\leq m+1$.
The new ratios are
\begin{equation}\label{eq:newratios}
 S_0=0,\quad
 S_j=F_t(R_j,R_{j-1})\ (1\leq j\leq m),\quad
 S_{m+1}=R_m+1/t.
\end{equation}
Extend the old ratios by $R_{-1}=0$ and $R_{m+1}=+\infty$.
For each $1\leq j\leq m$, apply \eqref{eq:preserve} with
\[
 (h,u,v,w)=(R_{j-2},R_{j-1},R_j,R_{j+1}).
\]
Convexity gives both comparison inequalities at interior indices.
At $j=1$, the lower comparison is $0\geq-R_1$.
At $j=m$, \eqref{eq:endpoint} gives precisely the last ratio in
\eqref{eq:newratios}. These endpoint conventions apply simultaneously
when $m=1$. We obtain $S_{j+1}-2S_j+S_{j-1}\geq0$ at every center.
Since $S_1>0=S_0$, the successive increments are all strictly positive.
This also provides the monotonicity needed for the next multiplication.

For a single factor, $R_0=0<R_1=1/x_1$, and there are no convexity
centers to check. Induction on the number of factors proves the lemma.
\end{proof}

\begin{lemma}\label{lem:HU}
Fix $1\leq r\leq m$ and write
\[
 A=e_r,\quad B=e_{r-1},\quad C=e_{r-2},\quad E=e_{r-3}.
\]
Then $A,B>0$, and
\begin{equation}\label{eq:HU}
 H=B^2-AC>0,\qquad U=B^2C-2AC^2+ABE\geq0.
\end{equation}
These sign conclusions persist when all four coefficients are multiplied
by a common positive constant.
\end{lemma}

\begin{proof}
For $r=1$, one has $B=1$, $C=E=0$, so $H=1$ and $U=0$.
For $r\geq2$, strict increase in Lemma~\ref{lem:ratios} gives
$B/A>C/B$, and hence $H>0$.
For $r=2$, convexity at the first center, multiplied by $AB$, gives
$U=B^2-2A\geq0$. For $r\geq3$, all of $A,B,C$ are positive and
\[
 \frac{U}{ABC}=R_r-2R_{r-1}+R_{r-2}\geq0.
\]
Finally, scaling multiplies $H$ by the square and $U$ by the cube
of the common factor.
\end{proof}

The variable domain in these lemmas matters. For a related use of
elementary-symmetric reduction, Zeng, Liu and Ratnavelu express order-five
permanent inequalities in symmetric functions of unimodular variables
and verify explicit positivity decompositions in exact arithmetic
\cite{ZengLiuRatnavelu2026}. That setting uses phase
constraints; Lemmas~\ref{lem:ratios}--\ref{lem:HU} here use positive real
variables and the ordered real ratios they provide.

\section{Column convexity and the finite transition}

Insertion of the largest element in a permutation gives the recurrence
\begin{equation}\label{eq:recurrence}
 c(n+1,k)=nc(n,k)+c(n,k-1).
\end{equation}
Indeed, the largest element is either a singleton cycle or can be inserted
after any of the $n$ existing elements. Induction then gives
\begin{equation}\label{eq:product}
 \sum_{k=0}^n c(n,k)z^k
 =z(z+1)\cdots(z+n-1)
 =(n-1)!z\prod_{j=1}^{n-1}(1+z/j),\quad n\geq1.
\end{equation}

\begin{proposition}\label{prop:column}
The ratios $q_k(n)$ satisfy \eqref{eq:qconvex}, with the strictness
asserted in Theorem~\ref{thm:main}.
\end{proposition}

\begin{proof}
Fix $n\geq k\geq2$ and put
\[
 (A,B,C,E)=(c(n,k),c(n,k-1),c(n,k-2),c(n,k-3)).
\]
By \eqref{eq:product}, these are $(n-1)!$ times the four coefficients
in Lemma~\ref{lem:HU}, with $r=k-1$, $m=n-1$, and variables
$1,1/2,\ldots,1/(n-1)$. Thus $A,B>0$, $H>0$, and $U\geq0$.
Two uses of \eqref{eq:recurrence} yield
\begin{align*}
 Q_1&=nA+B, & q_k(n+1)&=\frac{nB+C}{Q_1},\\
 Q_2&=n(n+1)A+(2n+1)B+C,
 &q_k(n+2)&=\frac{n(n+1)B+(2n+1)C+E}{Q_2}.
\end{align*}
Both $Q_1$ and $Q_2$ are positive. The numerator of the second difference
over the positive denominator $AQ_1Q_2$ is
\begin{align*}
 P={}&BQ_1Q_2-2A(nB+C)Q_2\\
     &+AQ_1\bigl(n(n+1)B+(2n+1)C+E\bigr)\\
   ={}&(nA+B)H+n(2B^3-3ABC+A^2E)+U.
\end{align*}
The identity
\[
 2B^3-3ABC+A^2E=\frac{2H^2+AU}{B}
\]
therefore gives the positive decomposition
\begin{equation}\label{eq:positive}
 P=(nA+B)H+\frac{2n}{B}H^2+
       \left(1+\frac{nA}{B}\right)U>0.
\end{equation}
This proves strict convexity for $k\geq2$ on the full domain $n\geq k$.
For $k=1$, $q_1(n)=0$ for every $n\geq1$, so its second difference
vanishes. No division by $B$ is used in that case.
\end{proof}

\begin{proof}[Proof of Theorem~\ref{thm:main}]
By Proposition~\ref{prop:column}, the sequence
\[
 d_j=q_k(j)-q_k(j+1),\qquad j\geq k,
\]
is nonincreasing. Some $d_j$ must be less than $1$. Otherwise, for every
integer $t\geq0$, telescoping would give
\[
 0\leq q_k(k+t)=q_k(k)-\sum_{j=k}^{k+t-1}d_j
       \leq q_k(k)-t,
\]
which fails for an integer $t>q_k(k)$. Hence the minimum $j_0$ in
\eqref{eq:threshold} exists, and $N_k=j_0+1\geq k+1$.

For $n\geq k+1$, the recurrence gives the exact identity
\begin{align}
 \frac{\D(n)}{c(n-1,k)c(n,k)}
 &=\frac{c(n,k)}{c(n-1,k)}-\frac{c(n+1,k)}{c(n,k)}\notag\\
 &=q_k(n-1)-q_k(n)-1=d_{n-1}-1.\label{eq:sign}
\end{align}
The denominator is positive. If $k+1\leq n<N_k$, minimality of $j_0$
gives $d_{n-1}\geq1$. If $n\geq N_k$, monotonicity gives
$d_{n-1}\leq d_{j_0}<1$. These are the two asserted signs on the
positive support. At the remaining indices,
$\D(n)=0$ for $1\leq n<k$ and $\D(k)=1$.
For $k=1$, one has $d_j=0$, giving $N_1=2$.
\end{proof}

\begin{remark}
The threshold in \eqref{eq:threshold} is the first negative center.
The theorem allows equalities before that center and requires no
classification of them. For the separate boundary $k=0$, the conventions
give $D_0(n)=0$ for every $n\geq1$.
\end{remark}

\section{Exact verification and assistance}

The proof uses finite products, rational identities, induction, and
telescoping. Its quantified conclusion does not depend on finite testing
or an asymptotic estimate. Symbolic exploration used SymPy~1.14.0
\cite{MeurerEtAl2017}. A separate standard-library Python implementation
checks the rational identity in \eqref{eq:preserve} and the identity obtained by multiplying
\eqref{eq:positive} by $B$ through sparse integer-polynomial expansion.
Both residual polynomials are zero.

The accompanying \texttt{verify.py} also constructs Stirling rows by
polynomial multiplication and cross-checks the permutation cycle-type
formula
\[
 c(n,k)=\sum_{\substack{m_1,\ldots,m_n\geq0\\
                  \sum_i i m_i=n,\ \sum_i m_i=k}}
           \frac{n!}{\prod_i i^{m_i}m_i!}
\]
through $n=22$. It checks $72\,000$ determinant signs for
$1\leq k\leq200$, $1\leq n\leq360$, and $52\,100$ second differences
for $1\leq k\leq200$, $k\leq n\leq360$ (using rows through $362$).
It also tests rational instances of the preservation inequality, including
negative comparison values and the terminal limit, and elementary-symmetric
ratios for all $791$ multisets of lengths $1$ through $7$ from
$\{1/10000,1/3,1,7,10000\}$. The script uses exact arithmetic and records
its coverage in \texttt{verification.json}. These finite checks detect
implementation and transcription errors; the preceding proof establishes
the infinite statement.

Generative AI assisted proof exploration, verification-code development,
manuscript preparation, and review in a separate model context. The model
review is not human peer review. The algebraic identities and all arguments
needed for the theorem are displayed above.

\bibliographystyle{plainnat}
\bibliography{references}
\end{document}
